\section{A calculus of pairs}
\label{sec:calculus}

This section makes the informal contract of \Cref{sec:overview} precise:
languages, behaviors, and projections (\S\ref{sec:languages}); pairs and
their commuting squares (\S\ref{sec:pairs}); and the pasting theorem that
lets squares compose into paths, together with the projection-compatibility
condition that pasting actually requires (\S\ref{sec:pasting}). Fidelity,
coverage, and the end-to-end guarantee build on this base in
\Cref{sec:fidelity}.

\subsection{Languages, behaviors, projections}
\label{sec:languages}

\begin{definition}[Language]
\label{def:language}
A \emph{language} $A$ is a triple $(\Prog{A}, F_A, \sem{\cdot}{A})$ where
$\Prog{A}$ is a decidable set of \emph{programs}; $F_A$ is a finite set of
named \emph{observable fields}, each $f \in F_A$ with a value domain
$V_f$; and $\sem{\cdot}{A} : \Prog{A} \rightharpoonup \Beh{A}$ is the
\emph{reference semantics}, a partial function into behaviors.
An \emph{observable state} is a total map
$\sigma \in \Obs{A} = \prod_{f \in F_A} V_f$, and a \emph{behavior} is a
finite sequence of observable states,
$\Beh{A} = \Obs{A}^{*}$, recorded \emph{post-step}: the $i$-th state is
the observable state after the $i$-th transition.
\end{definition}

The only admission criterion for a language is that $\sem{\cdot}{A}$ is
\emph{definable} --- an ISA manual, a logic's model theory, a reaction-network
semantics, and a molecular-graph notation all qualify. Nothing in this
section assumes languages are executable or even computational.

A program here is a \emph{closed} configuration: free inputs, when present,
are part of $p$ (a program together with an input valuation). Open programs
--- programs quantified over their inputs --- reappear in
\S\ref{sec:endtoend}, where a solver decides over all inputs and a witness
closes one.

\begin{definition}[Interpreter]
\label{def:interpreter}
An \emph{interpreter} for $A$ is a computable partial function
$I_A : \Prog{A} \rightharpoonup \Beh{A}$ that is \emph{pure}: identical
input yields byte-identical output, on every host. Outside its domain it
fails with a typed $\unsupported$ verdict naming the rejected construct;
$\dom(I_A)$ is $A$'s \emph{covered fragment}.
\end{definition}

\begin{assumption}[Interpreter adequacy]
\label{asm:adequacy}
$I_A(p) = \sem{p}{A}$ for all $p \in \dom(I_A)$.
\end{assumption}

Adequacy is an assumption, not a theorem, and we keep it visible: it is the
irreducible residue in every trusted-computing-base computation of
\S\ref{sec:tcb}. It is discharged \emph{empirically}, by differential
testing against an independently produced oracle for the same semantics
(the Sail-generated RISC-V simulator; the pinned CPython runtime;
\Cref{sec:instantiation}), and its failure modes are cross-checked by the
same machinery that checks translators.

\begin{definition}[Projection]
\label{def:projection}
A \emph{projection} over $A$ is a subset $\proj \subseteq F_A$. It lifts
to states by restriction, $\proj(\sigma) = \sigma\!\restriction_{\proj}$,
and to behaviors pointwise,
$\proj(\sigma_1 \cdots \sigma_n) = \proj(\sigma_1) \cdots \proj(\sigma_n)$
(in particular, projected behaviors of different lengths are unequal).
Write $b \eqproj{\proj} b'$ for $\proj(b) = \proj(b')$.
\end{definition}

\subsection{Pairs and the commuting square}
\label{sec:pairs}

\begin{definition}[Pair]
\label{def:pair}
A \emph{pair} $P : A \to B$ is a tuple $(A, B, T, \carry, \proj_P)$ where
$T : \Prog{A} \rightharpoonup \Prog{B}$ is the \emph{translator};
$\carry : \Beh{B} \rightharpoonup \Beh{A}$ is the
\emph{target-to-source interpreter}, which re-expresses a target behavior
as a source behavior; and $\proj_P \subseteq F_A$ is the pair's
\emph{declared projection} --- exactly the source observables it promises
to preserve. $T$ and $\carry$ are pure (as in \Cref{def:interpreter});
$I_A$ and $I_B$ are owned by the languages and shared by every pair that
touches them. The pair's domain $\dom(P)$ is the set of $p$ with
$p \in \dom(I_A) \cap \dom(T)$, $T(p) \in \dom(I_B)$, and
$I_B(T(p)) \in \dom(\carry)$.
\end{definition}

\begin{definition}[Faithfulness]
\label{def:faithful}
$P$ is \emph{faithful at} $p \in \dom(P)$ when the square commutes at $p$,
i.e.\ $\proj_P(I_A(p)) = \proj_P(\carry(I_B(T(p))))$:
\[
\begin{tikzcd}[column sep=large]
p \arrow[r, "T"] \arrow[d, "I_A"'] & T(p) \arrow[d, "I_B"] \\
I_A(p) & \carry(I_B(T(p))) \arrow[l, "\eqproj{\proj_P}"']
\end{tikzcd}
\]
$P$ is faithful on $C \subseteq \dom(P)$ when it is faithful at every
$p \in C$.
\end{definition}

Three remarks. First, everything in \Cref{def:faithful} is computable:
faithfulness at a given $p$ is a \emph{decidable, executable check} --- run
both sides, compare under $\proj_P$. We call this check the
\emph{square oracle}. Second, when the check fails it fails at a
\emph{first} position: a least step index $i$ and a field
$f \in \proj_P$ on which the two sides differ. The oracle reports
$(i, f)$ --- a translator (or interpreter, or carry-back) defect, localized
to a step and a named observable. Third, the projection is part of the
contract, not a tuning knob: $\proj_P$ states what ``preserves meaning''
promises for this pair, and $F_A \setminus \proj_P$ --- the pair's
\emph{loss} --- states what it does not.

The carry-back $\carry$ is what makes a target-side result \emph{mean}
something at the source: a solver counterexample or a target trace is
re-expressed as a source behavior. It is pair-specific because the
correspondence it encodes is pair-specific; it is also the reason the
square has an oracle of the same expressiveness as the translator it
checks, rather than a weaker abstraction.

\subsection{Pasting: when squares compose}
\label{sec:pasting}

Two pairs compose when the middle language is shared:
$P_1 = (A, B, T_1, \carry_1, \proj_1)$ and
$P_2 = (B, C, T_2, \carry_2, \proj_2)$ yield the candidate composite
\[
P_2 \circ P_1 \;=\; (A,\, C,\; T_2 \circ T_1,\; \carry_1 \circ \carry_2,\;
\proj)
\]
for a projection $\proj \subseteq \proj_1$ to be determined, with
$\dom(P_2 \circ P_1)$ the evident composite domain. Folklore says the
squares paste. They do not paste for free: $P_2$ guarantees agreement of
the two $B$-behaviors only \emph{up to} $\proj_2$, and $\carry_1$ consumes
whole $B$-behaviors. If $\carry_1$ reads a $B$-field that $P_2$ discards,
the outer rectangle can fail while both inner squares hold. The missing
condition is that $\carry_1$ must not look outside what $P_2$ preserves:

\begin{definition}[Support]
\label{def:support}
$\carry_1$ is \emph{$(\proj_2 \Rightarrow \proj)$-supported} when for all
$b, b' \in \dom(\carry_1)$:
$\proj_2(b) = \proj_2(b')$ implies
$\proj(\carry_1(b)) = \proj(\carry_1(b'))$, and moreover
$\dom(\carry_1)$ is $\proj_2$-closed on the behaviors compared
($b \in \dom(\carry_1)$ and $\proj_2(b)=\proj_2(b')$ imply
$b' \in \dom(\carry_1)$).
\end{definition}

\begin{theorem}[Pasting]
\label{thm:pasting}
Let $p \in \dom(P_2 \circ P_1)$, let $P_1$ be faithful at $p$, let $P_2$
be faithful at $T_1(p)$, let $\proj \subseteq \proj_1$, and let
$\carry_1$ be $(\proj_2 \Rightarrow \proj)$-supported. Then, writing
$r = T_2(T_1(p))$, the composite $P_2 \circ P_1$ is faithful at $p$ with
respect to $\proj$:
\[
\proj(I_A(p)) \;=\; \proj\bigl(\carry_1(\carry_2(I_C(r)))\bigr).
\]
\end{theorem}

\begin{proof}
Write $q = T_1(p)$. Then $\proj(I_A(p)) = \proj(\carry_1(I_B(q)))$ since
$\proj \subseteq \proj_1$ and $P_1$ is faithful at $p$. Faithfulness of
$P_2$ at $q$ gives
$\proj_2(I_B(q)) = \proj_2(\carry_2(I_C(r)))$, so the
support condition applies to the two $B$-behaviors and yields
$\proj(\carry_1(I_B(q))) = \proj(\carry_1(\carry_2(I_C(r))))$.
Chain the equalities. (Full details in \Cref{app:proofs}.)
\end{proof}

\paragraph{Composed keep and loss}
In the implementable special case --- $\carry_1$ \emph{fieldwise} and
step-aligned, each source field $f$ computed from a dependency set
$\mathit{deps}(f) \subseteq F_B$ of target fields --- the support
condition has a syntactic reading, and the largest $\proj$ satisfying
\Cref{thm:pasting} is
\[
\keep(P_2 \circ P_1) \;=\;
  \{\, f \in \proj_1 \mid \mathit{deps}(f) \subseteq \proj_2 \,\},
\]
a source field survives the path iff $P_1$ keeps it and it is computed
only from target fields $P_2$ keeps. Dually
$\lost(P_2 \circ P_1) = F_A \setminus \keep(P_2 \circ P_1)$: path loss is
the union of per-hop losses, each pulled back along the carry-backs. This
is the formal content of the platform rule that a path must \emph{declare
its cumulative loss}: the observables a destination can still speak about
are exactly those no hop discarded. (When $\carry_1$ changes step
granularity --- one source step spanning several target steps, as in
C-to-ISA --- the same statement holds with $\mathit{deps}$ taken over the
spanned window; \Cref{app:proofs}.)

\begin{corollary}[Localization]
\label{cor:localization}
Under the support condition, if the composite square fails at $p$ (both
sides defined), then $P_1$'s square fails at $p$ or $P_2$'s square fails
at $T_1(p)$. Inductively, along a route
$R = P_n \circ \cdots \circ P_1$ satisfying the pairwise support
conditions, a failing outer rectangle implies a failing inner square at
some hop $i$ --- and running the square oracle hop by hop finds the least
such $i$, then the least step and field within it.
\end{corollary}

\Cref{cor:localization} is the debugging story of the whole platform: a
divergence anywhere along a route is never a mystery about the route; it
is an indictment of one hop, one step, one observable. It is also what
disagreement between \emph{routes} will lean on in \S\ref{sec:branching}.

\paragraph{Routes}
The registry of languages and pairs forms a directed multigraph; a
\emph{route} $R : A \to Z$ is a path through it, with composed translator
$T_R$, composed carry-back $\carry_R$, composed projection $\keep(R)$, and
faithfulness given by iterating \Cref{thm:pasting}. Purity composes too,
and gives caching:

\begin{proposition}[Determinism and caching]
\label{prop:cache}
If every component function of every hop of $R$ is pure, then $T_R$ and
$\carry_R$ are pure, and a content-addressed cache keyed by the input
hash and the component versions returns, on a hit, exactly what
recomputation would return --- across the whole route. If some hop is impure, re-running the route on the same input and
diffing bytes at every hop (\emph{twice-and-diff}) detects it and
localizes it to the first hop whose bytes differ.
\end{proposition}
